\documentclass[a4paper,11pt]{article}
\usepackage[margin=1in]{geometry}
\usepackage[utf8]{inputenc}
\usepackage[T1]{fontenc}
\usepackage{amsmath}
\usepackage{booktabs}
\usepackage{hyperref}
\usepackage{enumitem}
\setlist{noitemsep}

\title{Labwork 2: CUDA Device Information}
\author{Nguyen Phong Chau}
\date{}

\begin{document}
\maketitle

\section*{Device Summary}

The CUDA runtime is available and correctly detected by Numba. The notebook output confirms that the system has one CUDA-capable device.

\begin{itemize}
    \item CUDA available: Yes
    \item CUDA devices detected: 1
\end{itemize}

\begin{table}[h]
\centering
\renewcommand{\arraystretch}{1.15}
\begin{tabular}{p{0.42\textwidth}p{0.42\textwidth}}
\toprule
Parameter & Value \\
\midrule
GPU name & NVIDIA GeForce GTX 1650 Ti with Max-Q Design \\
Multiprocessor count & 16 \\
Compute capability & 7.5 \\
Free memory & 3647.50 MB \\
Total memory & 3714.69 MB \\
\bottomrule
\end{tabular}
\caption{Detected CUDA device information}
\end{table}

The detected GPU is an NVIDIA GeForce GTX 1650 Ti with Max-Q Design, with 16 multiprocessors and compute capability 7.5. This device is suitable for CUDA programming and matches the runtime information reported by the notebook.

\end{document}
